\scanpage{23}
\mspage{10}

Le sous-groupe $\unsure{\Sigma|E}_\tau$ de $\Sigma$ image inverse du $\mathbb Z/2$ engendré par $\tau$, est le groupe cartographique ``\unsure{tri-angulé|triangulé} pondéré'' de Malgoire, ayant comme \deleted{\unsure{relation}} \inserted{gén. $\tau_0,\tau_1,\tau_\infty$} avec

\[
\left\{\begin{aligned}
\deleted{\tau_1\tau_\infty={}}l_0&=\tau_{\deleted{\unsure{0}}\infty}\tau_{\deleted{\unsure{\infty}}1}\\
\deleted{\tau_\infty\tau_0={}}l_1&=\tau_{\deleted{\unsure{1}}0}\tau_\infty\\
l_\infty&=\tau_1\tau_0
\end{aligned}\right. \tag*{(17)}
\]

\margin{\uline{NB} --}

$\Sigma_\tau$ peut être considéré comme le groupe ayant comme \inserted{gén. les $\tau_i$, avec les} seules relations (16) -- on a un hom signature évident $\tau_i\longrightarrow-1$ de $\Sigma_\tau$ dans $\pm1$, dont le noyau est $\pi$ engendré par les $l_i$ décrits par (17) \ldots.

\centerline{\rule{1.2cm}{0.4pt}}

% slanted
\margin{On a les relations \unsure{$\rho$}}

Revenons à une description plus \unsure{symétrique|symbolique} des autom. extérieurs à lacets de $\Sigma$. On désigne par $\widehat\Sigma$ le complété profini de $\Sigma$, extension de $\widehat{\unsure{\mathscr T|\Gamma}}_{0,3}=\mathscr T_{03}=\unsure{\mathfrak S_3|\mathcal S_3}\times\mathbb Z/2$ par $\widehat\pi$, et par $\widehat{\widehat\Sigma}$ le groupe des autom. \deleted{extérieurs} lacets de $\widehat\pi$, extension de $\widehat{\widehat\Gamma}=\mathrm{Aut\,ext}_{\mathrm{lac}}(\widehat\pi)$ par $\widehat\pi$

\diagram{\begin{tikzcd}[ampersand replacement=\&,column sep=small,row sep=0.2em]
\text{(18)} \& 1\arrow[r] \& \widehat\pi\arrow[r] \& \widehat{\widehat\Sigma}\arrow[r] \& \widehat{\widehat\Gamma}\arrow[r]\arrow[d,equal] \& 1\\
\& \& \& \& \mathrm{Aut\,ext}_{\mathrm{lac}}(\widehat\pi) \&
\end{tikzcd}}
% archive sheet number: 22
